% !TEX TS-program = pdflatex
% Example: the Executive Summary of a Politecnico di Milano master's
% thesis ("summary" format), using the modular academicthesis class +
% academicthesis-polimi. Mirrors the official Scuola 3I template
% (Executive_Summary).
% Finds the packages in ../core, ../languages, ../universities (any editor).
% academicthesis-repo-paths.tex -- lets showcases find the packages in this
% repository with ANY editor or engine (Texifier/Texpad, TeXShop, TeXstudio,
% plain pdflatex), not only with latexmk.
%
% \input it on the first line of a showcase, BEFORE \documentclass:
%
%   % academicthesis-repo-paths.tex -- lets showcases find the packages in this
% repository with ANY editor or engine (Texifier/Texpad, TeXShop, TeXstudio,
% plain pdflatex), not only with latexmk.
%
% \input it on the first line of a showcase, BEFORE \documentclass:
%
%   % academicthesis-repo-paths.tex -- lets showcases find the packages in this
% repository with ANY editor or engine (Texifier/Texpad, TeXShop, TeXstudio,
% plain pdflatex), not only with latexmk.
%
% \input it on the first line of a showcase, BEFORE \documentclass:
%
%   \input{academicthesis-repo-paths}
%   \documentclass[university=luiss]{academicthesis}
%
% It only extends LaTeX's own file search list (\input@path), which
% \documentclass, \RequirePackage, \input, \IfFileExists and
% \includegraphics all consult after TeX's normal search. The packages
% themselves stay path-free and are still loaded purely by name. Paths are
% relative to this showcases/ folder.
%
% \input@path does not search subfolders, so add one line per new
% university folder here (languages/ and core/ are flat).
\makeatletter
\providecommand*\input@path{}
\edef\input@path{%
  \input@path
  {../core/}%
  {../languages/}%
  {../universities/luiss/}%
  {../universities/polimi/}%
  {../universities/bocconi/}%
}
\makeatother
\endinput

%   \documentclass[university=luiss]{academicthesis}
%
% It only extends LaTeX's own file search list (\input@path), which
% \documentclass, \RequirePackage, \input, \IfFileExists and
% \includegraphics all consult after TeX's normal search. The packages
% themselves stay path-free and are still loaded purely by name. Paths are
% relative to this showcases/ folder.
%
% \input@path does not search subfolders, so add one line per new
% university folder here (languages/ and core/ are flat).
\makeatletter
\providecommand*\input@path{}
\edef\input@path{%
  \input@path
  {../core/}%
  {../languages/}%
  {../universities/luiss/}%
  {../universities/polimi/}%
  {../universities/bocconi/}%
}
\makeatother
\endinput

%   \documentclass[university=luiss]{academicthesis}
%
% It only extends LaTeX's own file search list (\input@path), which
% \documentclass, \RequirePackage, \input, \IfFileExists and
% \includegraphics all consult after TeX's normal search. The packages
% themselves stay path-free and are still loaded purely by name. Paths are
% relative to this showcases/ folder.
%
% \input@path does not search subfolders, so add one line per new
% university folder here (languages/ and core/ are flat).
\makeatletter
\providecommand*\input@path{}
\edef\input@path{%
  \input@path
  {../core/}%
  {../languages/}%
  {../universities/luiss/}%
  {../universities/polimi/}%
  {../universities/bocconi/}%
}
\makeatother
\endinput

\documentclass[university=polimi, master, format=summary, language=english]{academicthesis}

% The class only loads what the PoliMi layout needs; anything else is
% your choice. These are the ones the official template example uses.
\usepackage{amssymb}
\usepackage{bm}
\usepackage{float}    % [H] placement
\usepackage{colortbl} % \rowcolor{bluepoli!40}
\usepackage[colorlinks=true, linkcolor=black, citecolor=black, urlcolor=black]{hyperref}

\thesissetup{
  title         = {Title of the thesis},
  candidate     = {Name Surname},
  degree        = {Xxxxxxxxxxxx Engineering - Ingegneria Xxxxxxxxxxxx},
  supervisor    = {Prof. Name Surname},
  cosupervisor  = {Name Surname},
  academicyear  = {20XX-20XX},
}

%\addbibresource{example.bib}

\begin{document}

\frontpage

\section{Introduction}
\label{sec:introduction}
This document is an example of the PoliMi Executive Summary built with the
modular \texttt{academicthesis} class.

The Executive Summary is required only if the thesis has been assigned to
a reviewer (\textit{controrelatore}) for an independent evaluation of its
quality, scientific/technical relevance and original contribution.

\section{Guidelines}
The Executive Summary is a critical overview of your thesis with a focus
on the main achievements that have emerged from your research. It should
be organized in sections and paragraphs in order to better highlight the
major points of your work. Its length should range from four to six pages,
depending on the length of the thesis manuscript.

\section{Equations, Figures, Tables}
\subsection{Equations}
Maxwell's equations read:
\begin{align}
  \nabla\cdot \bm{D} & = \rho, \\
  \nabla \times \bm{E} + \frac{\partial \bm{B}}{\partial t} & = \bm{0}, \\
  \nabla\cdot \bm{B} & = 0, \\
  \nabla \times \bm{H} - \frac{\partial \bm{D}}{\partial t} & = \bm{J}.
\end{align}

\subsection{Figures}
\begin{figure}[H]
  \centering
  \includegraphics[width=0.6\columnwidth]{logopolimi3i}
  \caption{Caption of the Figure.}
  \label{fig:logo}
\end{figure}

\subsection{Tables}
\begin{table}[H]
  \centering
  \begin{tabular}{l c c c}
    \rowcolor{bluepoli!40}
    & \textbf{column1} & \textbf{column2} & \textbf{column3} \\
    \textbf{row1} & 1 & 2 & 3 \\
    \textbf{row2} & $\alpha$ & $\beta$ & $\gamma$ \\
  \end{tabular}
  \caption{Caption of the Table.}
  \label{table:example}
\end{table}

\section{Theorems and propositions}
\begin{theorem}
  Write here your theorem.
\end{theorem}

\begin{proposition}
  Write here your proposition.
\end{proposition}

\section{Conclusions}
A final section containing the main conclusions of your research.

\begin{acknowledgements}
  Here you might want to acknowledge someone.
\end{acknowledgements}

%\printbibliography

\end{document}
